\documentclass{report}
\usepackage{hyperref}
\usepackage{tcolorbox,float}
\usepackage{listings}

\makeatletter
\newfloat{info@box}{tbp}{loi}[section]% 1: Name of float environment. 2: Default placement (top, bottom, ...). 3: File extension if written to an aux-file (like toc, lof, lot, loa, ...). 4: Numbering within <section/subsection/...>.
\makeatother
\floatname{info@box}{Infobox}% Adapt caption.
\newenvironment{infobox}[1][]{% Create new environment using info@box and tcolorbox
   \begin{info@box}%
      \begin{tcolorbox}[colback=red!15!white,%                    background color
            colframe=red!75!black,%                               frame color
            title=Additional information\ifstrempty{#1}{}{: #1}.% title
         ]%
}{%
      \end{tcolorbox}%
   \end{info@box}%
}

\title{SoCLabs XiP AHB QSPI Configuration and User Manual}
\author{Daniel Newbrook, \href{http://www.soclabs.org}{SoC Labs}}
\begin{document}
\maketitle
\tableofcontents
\clearpage
\chapter{Introduction}
In SoC designs it is often quite valuable to have external non-volatile memory. Typically this will be 
where instruction code can be stored, and ideally we don't want to have to program this memory space after 
every power cycle. 

QSPI flash is fairly inexpensive and the protocol is fairly simple to implement. Whilst it isn't the fastest
protocol that you could use, it also doesn't require any real effort in the analog domain unlike for example SATA.

Excecute in place (XiP) is the method of making an external flash look like a memory mapped region to the processor.
Meaning the processor can directly run instructions from this address space. We add a small cache between the 
QSPI controller and the bus to reduce access time for repeat accesses. 

The design here uses an AHB interface for the main XiP interface, and an APB interface for access to internal 
registers of this IP, as well as to send commands over the QSPI interface.

\section{IP}
The SoCLabs AHB QSPI relies on IP's from certain vendors. In order to build you will need the following IP.
\subsection{Arm(R)}
Arm IP should be downloaded from Arm and placed into the recommended IP directories.
\begin{itemize}
    \item Corstone-101
    \item CG092 (flash cache)
\end{itemize}

\section{Integrating this IP}
There are 2 steps to integrating this IP in an SoC 
\begin{enumerate}
    \item Include the IP in the file list
    \item Instantiate the IP in the design
\end{enumerate}

For step 1. you make an envionment variable at your top level project pointing to the ahb qspi directory called 
SOCLABS\_AHB\_QSPI\_DIR (this will help with path references). You should then include the ahb\_QSPI\_SIM.flist to your 
behavioural flist, and ahb\_QSPI\_ASIC.flist to your asic flist. 

For step 2. you need an APB and AHB port to connect to, and then to bring out the QSPI wires to the top of your design. 
I think connecting these to simple tristate pads should be fine.

\chapter{Programmers Manual}

\section{Address Map}

The APB interface can access a number of registers. Outlined below are only the ones relating to the developed IP here,
for the CG092 register please see documentation from Arm.

\begin{table}[!h]
    \begin{center}
    \caption{Register Address Map}
    \label{address-map}
    \begin{tabular}{||c | c ||}
        \hline
        Register & Address \\ 
        \hline\hline
        Control Reg & 0x0000 \\
        \hline
        Status Reg & 0x0004 \\
        \hline
        SPI Commands & 0x0008 \\
        \hline
        SPI Address & 0x000C \\
        \hline
        Read Data 0 & 0x0010 \\
        \hline
        Read Data 1 & 0x0014 \\
        \hline
        Read Data 2 & 0x0018 \\
        \hline
        Read Data 3 & 0x001C \\
        \hline
        Write Data 0 & 0x0020 \\
        \hline
        Write Data 1 & 0x0024 \\
        \hline
        Write Data 2 & 0x0028 \\
        \hline
        Write Data 3 & 0x002C \\
        \hline
        AHB SPI Setup & 0x0030 \\
        \hline
        Clock Divider & 0x0034 \\
        \hline
        Peripheral ID 4-7 & 0x0FD0-0x0FDC \\
        \hline
        Peripheral ID 0-3 & 0x0FE0-0x0FEC \\
        \hline
        Component ID 0-3 & 0x0FF0-0x0FFC \\
        \hline
        CG092 Cache control & 0x1000-0x1FFF \\
        \hline
    \end{tabular}
\end{center}
\end{table}

\begin{table}[!h]
    \begin{center}
    \caption{Control Register Bits}
    \label{address-map}
    \begin{tabular}{||c | c ||}
        \hline
        Bits & Function \\ 
        \hline\hline
        0 & QSPI Quad IO Mode \\
        \hline 
        8 & XiP Active \\
        \hline
        16-23 & QSPI Mode Code \\
        \hline
        24 & Continous Read Active \\
        \hline 
        25 & No CMD active \\
        \hline
    \end{tabular}
\end{center}
\end{table}

\begin{table}[!h]
    \begin{center}
    \caption{Status Register Bits}
    \label{address-map}
    \begin{tabular}{||c | c ||}
        \hline
        Bits & Function \\
        \hline\hline
        0 & QSPI Busy (read-only) \\
        \hline
        8 & IRQ Clear (write 1 to clear QSPI finished interrupt) \\
        \hline
    \end{tabular}
\end{center}
\end{table}

\begin{table}[!h]
    \begin{center}
    \caption{SPI Commands Register Bits}
    \label{address-map}
    \begin{tabular}{||c | c ||}
        \hline
        Bits & Function \\ 
        \hline\hline
        0-7 & QSPI Command \\
        \hline 
        8 & QSPI Enable \\
        \hline
        9 & QSPI Read Enable \\
        \hline
        10 & QSPI Write Enable \\
        \hline 
        11 & QSPI Address enable \\
        \hline
        12-15 & Number Dummy cycles (N-1) \\
        \hline
        16-19 & Number R/W Bytes (N-1) \\
        \hline
    \end{tabular}
\end{center}
\end{table}

\begin{table}[!h]
    \begin{center}
    \caption{QSPI Address Register Bits}
    \label{address-map}
    \begin{tabular}{||c | c ||}
        \hline
        Bits & Function \\ 
        \hline\hline
        0-21 & QSPI Address \\
        \hline 
    \end{tabular}
\end{center}
\end{table}

\begin{table}[!h]
    \begin{center}
    \caption{AHB SPI Setup Register Bits (0x0030)}
    \begin{tabular}{||c | c ||}
        \hline
        Bits & Function \\
        \hline\hline
        0-7 & AHB QSPI Command (opcode for XiP reads) \\
        \hline
        12-15 & AHB Dummy Cycles (N-1) \\
        \hline
    \end{tabular}
\end{center}
\end{table}

\begin{table}[!h]
    \begin{center}
    \caption{Clock Divider Register Bits (0x0034)}
    \begin{tabular}{||c | c ||}
        \hline
        Bits & Function \\
        \hline\hline
        0-4 & QSPI Clock Divider (SCLK = HCLK / 2N, default=1) \\
        \hline
    \end{tabular}
\end{center}
\end{table}

\section{Interrupts}
The QSPI controller provides a single interrupt output (\texttt{IRQ\_QSPI\_FINISHED}) which is asserted on the falling edge of the QSPI busy signal (i.e. when a QSPI transaction completes). To clear the interrupt, write a 1 to bit 8 of the Status Register (0x0004).

\section{Setting XiP Mode}
Typical flow for entering XiP
\begin{enumerate}
    \item Reset the QSPI flash
    \item Enable QSPI Mode
    \item Set Continous Read mode 
    \item Set XiP/AHB Mode
    \item Configure Cache 
\end{enumerate}

\section{Write to Flash}
Typical flow for writing to flash 
\begin{enumerate}
    \item Reset the QSPI Flash
    \item Enable QSPI mode
    \item Clear write protection
    \item Set write enable 
    \item Write to APB registers
    \item Write to flash
    \item Poll flash status until complete
\end{enumerate}

\end{document}